\section{Related Work}

Existing information retrieval test collections for evaluating retrieval and screening in systematic reviewing have predominantly focused on the biomedical domain. While these datasets differ in scale, annotation depth, and supported tasks, they share a common goal: providing reusable test collections for studying search (i.e., retrieving literature), screening (i.e., determining if literature should be included in the review or not), and study prioritization (i.e., ranking the literature to assist with the screening phase) in systematic reviews. Table~\ref{table-sr-datasets} summarizes the most commonly used resources and contrasts them with Webis-SR4ALL-26.

We first discuss typical test collections that are commonly used in systematic review automation tasks. The SIGIR~2017 SysRev Query Collection~\cite{scells:2017} comprises Cochrane reviews published between~2014 and~2016. The original Boolean search strategies were translated into executable Elasticsearch queries. Included and excluded studies were mapped to PubMed identifiers. The Seed Studies Collection~\cite{wang:2022} targets seed-driven search strategies in systematic reviewing; i.e., the creation of Boolean queries driven by knowledge about a priori studies that are likely to be included in the review. It provides rich review-process artifacts, including original PubMed Boolean queries, seed studies used to assist in the formulation of the Boolean queries, and lists of included studies. In addition, the collection includes references obtained through automated citation chasing tools such as CitationChaser and SpiderCite. The CLEF TAR collections~\cite{kanoulas:2017,kanoulas:2018,kanoulas:2019} were released as part of the CLEF lab on technology-assisted reviews in empirical medicine. Built on Cochrane reviews, it contains expert-formulated queries, retrieved records, and relevance assessments spanning intervention, prognosis, qualitative studies, and diagnostic test accuracy review types. The corpus primarily supports querying and screening prioritization tasks. 

Apart from the typical retrieval and screening tasks, \citet{alharbi:2019} proposed a dataset for performing these tasks on systematic review updates, where there are two versions of the review and the task is to assist in helping identify relevant studies for the updated version. The dataset from~\citet{hannousse:2021} is used exclusively for screening and is the only other non-biomedical systematic review collection we identified in our review of the literature. The CSMeD dataset~\cite{kusa:2023} consolidated prior biomedical and computer science screening datasets~\cite{cohen:2006,wallace:2010,howard:2016,scells:2017,kanoulas:2017,kanoulas:2018,kanoulas:2019,alharbi:2019,hannousse:2021} into a unified meta-collection of 325~systematic reviews. By harmonizing nine existing resources and providing standardized access via the BigBio framework, CSMeD enables more comparable evaluations of citation screening methods. The dataset further enriches a subset of Cochrane reviews by including eligibility criteria and search strategies extracted from review protocols, as well as CSMeD-FT, a full-text screening collection with time-based splits to prevent data leakage. Finally, the most recent dataset for systematic review automation we identified was from \citet{wang:2025b}. This dataset was used primarily as training data for Boolean query formulation through reinforcement learning. While it is by far the largest dataset compared to all others discussed above, it provides fewer review-process artifacts and is not as useful for many other tasks for automating systematic reviews.

While these resources have been instrumental in advancing retrieval and screening methods for biomedical systematic reviews, their scope usually remains limited to a single scientific domain and a small number of curated review topics. In contrast, Webis-SR4ALL-26 is desigend as a large-scale, cross-disciplinary corpus covering hundreds of thousands of systematic reviews from a variety of scientific fields. It enables research into systematic reviews at a scale previously not possible, and across tasks like retrieval and screening. With further extraction and annotation effort, it may also support future tasks like syntehsizing systematic reviews.